\chapter{全纯函数}
\label{chap:holomorphic-functions-solutions}

\begin{solution}{complex-spherical-convergence}
按倒数约定，直接代入弦距离公式可得
\[
 \chi(1/z,0)=\chi(z,\infty)\qquad(z\in\widehat{\C}).
\]
这在 \(z=0,\infty\) 时分别是两边等于一、两边等于零。因此所述两种
收敛等价，也包括序列反复取无穷值的情形。对于有限 \(z\)，又有

\[
 \chi(z,\infty)<\varepsilon
 \iff \frac1{\sqrt{1+|z|^2}}<\varepsilon
 \iff |z|>\sqrt{\varepsilon^{-2}-1}
\]

对 \(0<\varepsilon<1\) 成立。加上球心 \(\infty\) 后，弦距离球恰等于
正文中的一个无穷点基本邻域；随着 \(\varepsilon\downarrow0\)，圆外半径
趋于无穷，两种邻域定义一致。
\end{solution}

\begin{solution}{complex-limit-components}
若 \(|f(z)-L|<\varepsilon\)，则

\[
 |\operatorname{Re}f(z)-\operatorname{Re}L|
 \le|f(z)-L|,\qquad
 |\operatorname{Im}f(z)-\operatorname{Im}L|
 \le|f(z)-L|,
\]

故两个分量分别收敛。反之，若两分量误差都小于
\(\varepsilon/\sqrt2\)，则

\[
 |f(z)-L|^2
 =|\operatorname{Re}(f-L)|^2+|\operatorname{Im}(f-L)|^2
 <\varepsilon^2.
\]

对乘法，写
\[
 f(z)g(z)-LM=f(z)(g(z)-M)+M(f(z)-L).
\]
收敛函数在充分小的删心邻域有界，故右端两项分别趋零。
\end{solution}

\begin{solution}{cr-point-counterexample}
沿坐标轴计算：
\[
 \frac{u(h,0)-u(0,0)}h=1,\quad
 \frac{u(0,h)-u(0,0)}h=0,
\]
\[
 \frac{v(h,0)-v(0,0)}h=0,\quad
 \frac{v(0,h)-v(0,0)}h=1.
\]
所以 \(u_x=v_y=1\)、\(u_y=-v_x=0\)。沿 \(h=t\in\R\)，复差商为
\(1\)。沿 \(h=t(1+i)\)，有
\[
 f(t+it)=\frac t2+i\frac t2=\frac12t(1+i),
\]
故差商为 \(1/2\)。两个方向极限不同，复导数不存在。这也说明在一点仅核对
偏导和 Cauchy--Riemann 方程不能代替实 Fréchet 可微性。
\end{solution}

\begin{solution}{complex-linear-real-matrix}
在实基 \(1,i\) 下设 \(L(1)=a+ib\)。若 \(L\) 复线性，则
\(L(i)=iL(1)=-b+ia\)，故两列分别为 \((a,b)^{\mathsf T}\) 与
\((-b,a)^{\mathsf T}\)。反之，具有该矩阵的实线性映射满足
\[
 L(x+iy)=x(a+ib)+y(-b+ia)=(a+ib)(x+iy),
\]
所以 \(L(\lambda z)=\lambda L(z)\) 对所有 \(\lambda,z\in\C\) 成立。
\end{solution}

\begin{solution}{continuous-cr-sufficiency}
固定 \((x_0,y_0)\)。连续偏导给出
\[
 \begin{aligned}
 u(x_0+h_1,y_0+h_2)-u(x_0,y_0)
 &=u_x(x_0,y_0)h_1+u_y(x_0,y_0)h_2+o(|h|),\\
 v(x_0+h_1,y_0+h_2)-v(x_0,y_0)
 &=v_x(x_0,y_0)h_1+v_y(x_0,y_0)h_2+o(|h|).
 \end{aligned}
\]
可用先沿水平、再沿竖直线段的一元中值定理证明这两个展开，连续性保证两处
偏导可换成基点偏导而只产生 \(o(|h|)\)。令
\(a=u_x=v_y\)、\(b=v_x=-u_y\)，合并两式得到
\[
 f(z_0+h)-f(z_0)=(a+ib)h+o(|h|).
\]
除以 \(h\ne0\) 后余项模为 \(o(1)\)，故复导数存在且等于 \(a+ib\)。
\end{solution}

\begin{solution}{holomorphic-real-line-image}
设像包含在直线 \(w_0+e^{i\theta}\R\) 中。函数
\[
 g(z)=e^{-i\theta}(f(z)-w_0)
\]
全纯且取实值。写 \(g=u+iv\)，则 \(v\equiv0\)。Cauchy--Riemann 方程给
\(u_x=v_y=0\)、\(u_y=-v_x=0\)，所以 \(g'=0\)。域连通，故 \(g\) 常值，
从而 \(f\) 常值。
\end{solution}

\begin{solution}{holomorphic-constant-modulus}
若恒定模为零，显然 \(f\equiv0\)。设 \(|f|=c>0\)，写 \(f=u+iv\)。
由 \(u^2+v^2=c^2\) 分别对 \(x,y\) 求导，并用 Cauchy--Riemann 方程，得到
\[
 \begin{pmatrix}u&v\\-v&u\end{pmatrix}
 \binom{u_x}{v_x}=0.
\]
矩阵行列式为 \(u^2+v^2=c^2>0\)，故 \(u_x=v_x=0\)，CR 方程再给
\(u_y=v_y=0\)。于是 \(f'=0\)，连通性说明 \(f\) 为常数。
\end{solution}

\begin{solution}{complex-power-series-derivatives}
内闭一致收敛和逐项求导允许反复求导：
\[
 f^{(k)}(z)=
 \sum_{n=k}^{\infty}n(n-1)\cdots(n-k+1)a_n(z-z_0)^{n-k}.
\]
取 \(z=z_0\) 时只有 \(n=k\) 项留下，故 \(f^{(k)}(z_0)=k!a_k\)。
若另一个系数列 \(b_n\) 在同一中心的某邻域表示 \(f\)，同样计算给
\(k!a_k=f^{(k)}(z_0)=k!b_k\)，所以 \(a_k=b_k\) 对所有 \(k\) 成立。
\end{solution}

\begin{solution}{complex-exponential-fibres}
由 \(e^{x+iy}=e^x(\cos y+i\sin y)\)，若 \(e^z=e^w\)，比较模得到
\(e^{\operatorname{Re}z}=e^{\operatorname{Re}w}\)，故实部相等。再比较单位
复数得到
\[
 e^{i(\operatorname{Im}z-\operatorname{Im}w)}=1,
\]
即虚部之差属于 \(2\pi\Z\)。逆向由周期性立即成立。因此若
\(w\ne0\)，选取任一辐角 \(\theta\)，则
\[
 \exp^{-1}(\{w\})
 =\{\log|w|+i(\theta+2k\pi):k\in\Z\}.
\]
\(w=0\) 没有逆像。
\end{solution}

\begin{solution}{principal-logarithm-derivative}
在裂平面上令
\[
 u=\frac12\log(x^2+y^2),\qquad
 v=\operatorname{Arg}(x+iy)\in(-\pi,\pi).
\]
局部选取连续辐角后，
\[
 u_x=\frac{x}{x^2+y^2},\quad u_y=\frac{y}{x^2+y^2},\quad
 v_x=-\frac{y}{x^2+y^2},\quad v_y=\frac{x}{x^2+y^2}.
\]
偏导连续且满足 \(u_x=v_y\)、\(u_y=-v_x\)，故主值对数全纯。其导数为
\[
 u_x+iv_x=\frac{x-iy}{x^2+y^2}=\frac1{x+iy}.
\]
切口使辐角能单值连续选择；跨过负实轴时上下极限相差 \(2\pi\)。
\end{solution}

\begin{solution}{wirtinger-real-linear-decomposition}
令
\[
 a=\frac12\bigl(L(1)-iL(i)\bigr),\qquad
 b=\frac12\bigl(L(1)+iL(i)\bigr).
\]
则 \(a+b=L(1)\)、\(i(a-b)=L(i)\)。对 \(h=x+iy\)，实线性给
\[
 L(h)=xL(1)+yL(i)=a(x+iy)+b(x-iy)=ah+b\bar h.
\]
取 \(h=1,i\) 可反解出 \(a,b\)，所以表示唯一。对 \(L=Df(z)\)，由
\[
 Df(z)(1)=f_x,\qquad Df(z)(i)=f_y
\]
代入上式得到 \(a=(f_x-if_y)/2=f_z\)、
\(b=(f_x+if_y)/2=f_{\bar z}\)，即正文公式。
\end{solution}

\begin{solution}{polar-cauchy-riemann}
由 \(x=r\cos\theta\)、\(y=r\sin\theta\)，链式法则给
\[
 u_r=u_x\cos\theta+u_y\sin\theta,\qquad
 \frac1r u_\theta=-u_x\sin\theta+u_y\cos\theta,
\]
对 \(v\) 也有同样公式。把 \(u_x=v_y\)、\(u_y=-v_x\) 代入，整理得
\[
 u_r=\frac1r v_\theta,\qquad v_r=-\frac1r u_\theta.
\]
反向由旋转矩阵可逆同样成立。对 \(z^n=r^ne^{in\theta}\)，
\(u=r^n\cos n\theta\)、\(v=r^n\sin n\theta\)，直接求导即满足两式；
导数为 \(nz^{n-1}\)。负整数时结论在 \(r>0\) 的区域成立。
\end{solution}

\begin{solution}{local-holomorphic-inverse}
实 Jacobian 行列式为 \(|f'(z_0)|^2>0\)。实逆函数定理给邻域
\(U\ni z_0\)、\(V\ni f(z_0)\)，使 \(f:U\to V\) 为 \(C^1\) 微分同胚。
令 \(g=f^{-1}\)。链式法则给
\[
 Dg(f(z))=[Df(z)]^{-1}.
\]
\(Df(z)\) 是乘以 \(f'(z)\) 的复线性映射；缩小 \(U\) 后
\(f'(z)\ne0\)，其逆是乘以 \(1/f'(z)\)，仍复线性。实微分判据说明
\(g\) 全纯，并有
\[
 g'(w)=\frac1{f'(g(w))}.
\]
\end{solution}
